  
\addcontentsline{toc}{chapter}{\twkkai{对话录·其五}}  
\markboth{\twkkai{对话录·其五}}{}  
\pagecolor{black}
\color{white}
\quad\quad \twkkai{\lettrine[,lraise=0.2,lhang=0.5]{\textbf{猫}}{：你去哪了啊？打了个盹就不见你影了喵。}

狄利克雷：还是要由我来摆平这个烂摊子啊……

猫：你……在责备我办事不力？

狄利克雷：倒没有这个意思。我只是强调，有些事还是要亲力亲为而已。

（猫低垂着头，悄悄地跃过台阶，伸起了懒腰；同时用爪子摩挲着墙\marginnote{
\centering
{\Large ♪}

\vspace{0.4em}

\qraudio{quartet1}{\quartetSound}{https://raw.githubusercontent.com/kelvinlamresearchmath-web/finding-dirichlet-site/master/quartet1.mp3}

\vspace{0.3em}
}[0.5cm]
{\small ♪}。）

\noindent\adjustbox{center}{\includegraphics[scale=0.40]{quartet1_white.pdf}}


狄利克雷：我知道你很惆怅——但你最好把你的爪子从墙上拿下来。

猫：喵……你不觉得，我们之间新的一轮讨论，就要开始了吗？

狄利克雷：噢，是吗？那我觉得，我们倒应该好好聊聊奖章的问题了。

猫：你这不就是责备我办事不力了嘛……

狄利克雷：看来你也知道自己丢失奖章很久了，是吗？

猫：（支吾着）好吧；我已经尽全力了，那两个使者把奖章传来传去的……

狄利克雷：你竟然还让他们把奖章传来传去？你知道这对于我们天庭有多危险吗？

猫：我知错了……

守门人：……163……

狄利克雷：这守门人，让他完成的任务他不去做，反而在这给我念叨我给他的工号。

猫：163？这可是一个很神奇的数字啊！

狄利克雷：你又打岔了。

猫：不不！刚刚我刮响的那一段门德尔松Op.81第三乐章（Capriccio）的引子，和这个或许有微弱的联系呢！

狄利克雷：好，好——讲吧。

猫：这又回到了我们一开始——也就是很久以前——聊到的动机（\textit{motif}）了。你看前几个音符，重复着：\kern-15pt\lily[hpadding=0pt]{
    \override Staff.StaffSymbol.color       = #white
    \override Staff.Clef.color              = #white
    \override Staff.TimeSignature.color     = #white
    \override Staff.KeySignature.color      = #white
    \override Staff.BarLine.color           = #white
    \override Staff.LedgerLineSpanner.color = #white
    \override NoteHead.color                = #white
    \override Stem.color                    = #white
    \override Beam.color                    = #white
    \override Flag.color                    = #white
    \override Rest.color                    = #white
    \override Accidental.color              = #white
    \override Dots.color                    = #white
    \override Slur.color                    = #white
    \override Tie.color                     = #white
    \override Script.color                  = #white
    \override DynamicText.color             = #white
    \override Hairpin.color                 = #white
    \override TextScript.color              = #white
    \override TupletBracket.color           = #white
    \override TupletNumber.color            = #white
  \clef treble
  \key g \major
  \time 12/8
\relative c' {
\partial 2. r8 b( e g fis e) r8 b( e g fis e) r8 c( e a g fis) 
} \bar "||"
}\kern5pt 。

狄利克雷：这段让人感到很宁静呢。

猫：是啊——它的速度标记是\textit{Andante con moto}；不过很快我就会给你介绍这同一首曲子后面突然变得热闹的部分——但不是现在讲喵。

狄利克雷：好。那现在这段和163的联系在哪呢？

猫：不不——这个时候还没轮到163这位主角呢！是另一位潜伏者和我刚才反复哼唱的重复动机有联系。

狄利克雷：具体是……？

守门人：……\(
\boxed{
j(\tau)
=
\frac{E_4(\tau)^3}{\Delta(\tau)}
}
\)
……它不会随着模群的变换而改变——
\(
\tau \mapsto \frac{a\tau+b}{c\tau+d}
\)……无论怎么变换——\(
j\!\left(\frac{a\tau+b}{c\tau+d}\right) = j(\tau)
\)……

猫：很好，很好喵！守门人精确地帮我回答了你的问题——那个潜伏者，是\(j(\tau)\)！

狄利克雷：那它们怎么会相似呢？

猫：请允许我再重复一遍吧——\kern-15pt\lily[hpadding=0pt]{
    \override Staff.StaffSymbol.color       = #white
    \override Staff.Clef.color              = #white
    \override Staff.TimeSignature.color     = #white
    \override Staff.KeySignature.color      = #white
    \override Staff.BarLine.color           = #white
    \override Staff.LedgerLineSpanner.color = #white
    \override NoteHead.color                = #white
    \override Stem.color                    = #white
    \override Beam.color                    = #white
    \override Rest.color                    = #white
    \override Accidental.color              = #white
    \override Dots.color                    = #white
    \override Slur.color                    = #white
    \override Tie.color                     = #white
    \override Script.color                  = #white
    \override DynamicText.color             = #white
    \override Hairpin.color                 = #white
    \override TextScript.color              = #white
    \override TupletBracket.color           = #white
    \override TupletNumber.color            = #white
  \clef treble
  \key g \major
  \time 12/8
\relative c' {
\partial 2. r8 b( e g fis e) r8 b( e g fis e) r8 c( e a g fis) 
} \bar "||"
}\kern5pt 。这段里，每一组动机由一个八分休止符和五个八分音符构成，它们的变化就像\(\tau\)这个自变量变了之后\(j(\tau)\)发生的相应的变化。但是万变不离其宗——\(j(\tau)\)的变化是有规律可循的——像守门人说的：它不会随着模群的变换而改变——\(\tau \mapsto \frac{a\tau+b}{c\tau+d}\)。无论怎么变换——\(j\!\left(\frac{a\tau+b}{c\tau+d}\right) = j(\tau)\)！

狄利克雷：原来如此！那这不是绕回到我们一开始的……变中的不变性？

猫：哈！的确，但这次变了，163登场了。

守门人：……我写下来吧：
\(
\tau_{163}
=
\frac{1+\sqrt{-163}}2
\)
对应：
\(
j\!\left(\frac{1+\sqrt{-163}}2\right)
\\=
-262537412640768000
=
-640320^3
\)……由于类数为 $1$：
\(
j(\tau)\in\mathbb Z.
\)

狄利克雷：守门人怎么了？怎么开始乱报数字……

猫：不不！这就是163开始发挥作用了——更准确地说，是\(\tau_{163}\)开始主宰整个舞台。你没发现，刚才他报的那一串数字，是一个整数吗？

狄利克雷：等等……还真是。

猫：是吧，这就是类数的威力了——163是最大的Heegner数呢！嗯，这下就圆满收尾了。伴随着悠扬的主旋律，163和\(j(\tau)\)齐齐地走上了舞台。

（猫哼起了小调：\kern-15pt\lily[hpadding=0pt]{
    \override Staff.StaffSymbol.color       = #white
    \override Staff.Clef.color              = #white
    \override Staff.TimeSignature.color     = #white
    \override Staff.KeySignature.color      = #white
    \override Staff.BarLine.color           = #white
    \override Staff.LedgerLineSpanner.color = #white
    \override NoteHead.color                = #white
    \override Stem.color                    = #white
    \override Beam.color                    = #white
    \override Flag.color                    = #white
    \override Rest.color                    = #white
    \override Accidental.color              = #white
    \override Dots.color                    = #white
    \override Slur.color                    = #white
    \override Tie.color                     = #white
    \override Script.color                  = #white
    \override DynamicText.color             = #white
    \override Hairpin.color                 = #white
    \override TextScript.color              = #white
    \override TupletBracket.color           = #white
    \override TupletNumber.color            = #white
  \clef treble
  \key g \major
  \time 12/8
\relative c' {
g''4.( fis4) dis8 fis4( e8) d4( b8) d4( c8) a4( fis8) fis4.( fis8)( g a)
} \bar "||"
}\kern5pt 。）

\rule{0pt}{20pt}

\noindent\adjustbox{center}{\includegraphics[scale=1]{trebleclef.preview.pdf}}


}













